\hypertarget{chapter-24-the-geometry-of-logical-regimes}{%
\chapter*{Chapter 24 --- The Geometry of Logical Regimes}
\addcontentsline{toc}{chapter}{Chapter 24 --- The Geometry of Logical Regimes}
\markboth{Chapter 24 --- The Geometry of Logical Regimes}{Chapter 24 --- The Geometry of Logical Regimes}\label{chapter-24-the-geometry-of-logical-regimes}}

Chapters 20 through 23 built three things over each proposition and each
cover: a support configuration (Chapter 20), a coherence defect on that
configuration (Chapter 21), and a repair cost that depends on both at
once (Chapter 23), with Chapter 22 establishing that the first two vary
independently. What hasn't been asked yet is what kind of \emph{space} all of
this lives in --- and whether calling it a geometry, rather than merely a
list of cases, is earned or aspirational. This chapter argues it's
earned, in one precise sense, and proposes the rest as a program rather
than claiming more than has been built.

\hypertarget{whats-actually-rigorous-t-f-b-n-as-fibers-not-points}{%
\section{1. What's actually rigorous: T, F, B, N as fibers, not points}\label{whats-actually-rigorous-t-f-b-n-as-fibers-not-points}}

Fix a cover \(\{U_i\}\) and a proposition \(p\). Define the \textbf{support-coherence
space} for \(p\),
\[\mathcal{S}(p) = \Big\{ \big(e^+_i, e^-_i\big)_i \ \Big|\ e^+_i \in \mathcal{E}^+(U_i),\ e^-_i \in \mathcal{E}^-(U_i) \Big\},\]
the set of all possible support assignments across the cover --- the full
data Chapter 20 works with, before anything is merged or compared. There
is a projection
\[\pi : \mathcal{S}(p) \to \mathbb{F}, \qquad \pi\big((e^+_i,e^-_i)_i\big) = \Sigma(p) = \bigvee_i q(e^+_i,e^-_i),\]
Chapter 22's merge. Chapter 22's central result, restated in this
chapter's terms, is that \(\pi^{-1}(T)\) is not a point --- it contains both
case (a) (coherent) and case (c) (incoherent), which are different
elements of \(\mathcal{S}(p)\) that happen to share an image under \(\pi\).
The same holds for every fiber. This is what it means, precisely rather
than suggestively, to say that \(T\), \(F\), \(B\), \(N\) are \textbf{coarse
projections of regions in a richer space}, not native points of the
space itself. Nothing here needed new machinery --- it's Chapter 22's
result, restated as a fact about a map between two spaces rather than as
a fact about two separately-computed quantities. That restatement is the
whole of what's rigorous in this chapter; everything past this section is
proposed rather than proven.

\hypertarget{why-a-taxonomy-isnt-yet-a-geometry}{%
\section{2. Why a taxonomy isn't yet a geometry}\label{why-a-taxonomy-isnt-yet-a-geometry}}

A taxonomy sorts \(\mathcal{S}(p)\) into labeled bins --- clean glut,
confounded glut, coherent-\(T\), incoherent-\(T\), and so on, exactly as
Chapter 22's six-cell table does. A geometry needs more: a notion of which
configurations are \emph{near} which others, and what it costs to move from
one to another. Without that, ``the space of logical regimes'' is a
figure of speech for a partition, and the chapter's title would be false
advertising. Two candidate structures are worth taking seriously, and
this is where TARTAN stops being a correspondence-matrix entry from an
earlier chapter and becomes load-bearing.

\hypertarget{neighborhoods-from-refinement-the-proposed-topology}{%
\section{3. Neighborhoods from refinement --- the proposed topology}\label{neighborhoods-from-refinement-the-proposed-topology}}

\(\mathcal{S}(p)\) as defined above is relative to one fixed cover
\(\{U_i\}\). TARTAN's recursive tiling gives a \emph{family} of covers, ordered
by refinement --- each tile decomposable into finer sub-tiles, with the
TARTAN Admissibility Theorem guaranteeing that global admissibility
reduces to tile-local admissibility at every level of that refinement,
not just at one chosen resolution. This suggests treating \(\mathcal{S}(p)\)
not as fixed to one cover but as assembled from the whole refinement
diagram --- coarser covers give coarser support-coherence spaces, finer
covers give finer ones, and the refinement maps between them are exactly
what a topology on the assembled space would need to be built from:
two configurations count as \emph{near} each other if they agree at some
coarse level of tiling and only differ once the tiling is refined past
some threshold.

This is proposed, not built. What would need proving: that refining a
cover is compatible with the coherence defect in the right way (does
refining a tile ever \emph{manufacture} an apparent coherence defect that
wasn't real at the coarser resolution, or only ever reveal genuine ones
that coarseness had hidden?), and that the resulting limit or colimit
construction is well-behaved enough to actually deserve the name
topology rather than merely a directed system of unrelated spaces. Both
are real technical questions, not merely items to gesture past.

\textbf{Running Part IV's test against this proposal.} Chapter 15 posed a
specific question for whoever built this topology to answer, rather than
leaving the relationship between dispersion obstruction and this
chapter's program unresolved indefinitely: does the topology sketched
above contain enough structure to measure \emph{persistence} under
refinement, or only \emph{neighborhoods}? As stated, it's the latter --- ``near''
here means agreeing at some coarse level and differing only past a
refinement threshold, which is a qualitative nearness relation with no
attached magnitude. Nothing in §3 as written could represent an
accumulation of small displacements the way Chapter 15's
\(\sum_i \epsilon_i\) does; there's no metric, uniformity, or filtration
here, only a directed system of covers. So the answer to Chapter 15's
test is settled, at least for this proposal in its current form:
\textbf{dispersion obstruction and this section's topology are not
duplicates.} This section, if built, would supply the qualitative
geometry --- which configurations are near which, at what refinement
depth. Dispersion needs the further scale structure Chapter 15
introduced on top of that, not in place of it. Whether this topology
should be \emph{extended} with that structure, or whether dispersion is
better left as a separate functional layered over a topology that stays
qualitative, is a design choice for whoever builds this section
properly --- this note only settles that the two aren't secretly the same
thing.

\hypertarget{distance-from-repair-cost-the-proposed-metric}{%
\section{4. Distance from repair cost --- the proposed metric}\label{distance-from-repair-cost-the-proposed-metric}}

Chapter 23's functional \(J(a) = E(a) + \lambda R(a) + \mu C(a)\) was built
to price an action's dependence on a given support-coherence
configuration, not to measure distance between two configurations
directly. But the ingredients are close enough to a metric to name the
gap precisely: for two configurations \(s, s' \in \mathcal{S}(p)\), define
a candidate transition cost as the minimal repair expenditure
(\(\lambda R + \mu C\) terms, stripped of the action-specific \(E\)) needed
to move from one to the other --- turning a confounded glut into a clean
one, or eliminating a coherence defect entirely, has some cost, and nearby
configurations should be the ones reachable cheaply.

This is also proposed rather than built, and for a specific, nameable
reason: \(J\) as defined in Chapter 23 is a cost \emph{of an action}, relative to
what that action needs from \(p\). A genuine distance on \(\mathcal{S}(p)\)
would need a notion of repair cost that doesn't route through any
particular action --- something closer to \(\inf_a\) over all actions, or a
repair cost defined directly on configurations rather than borrowed from
\(C(a,p)\). Whether that infimum is well-behaved, and whether it satisfies
anything like a triangle inequality, is exactly the kind of question this
chapter can pose but not yet answer.

\hypertarget{what-the-geometry-would-buy-if-built}{%
\section{5. What the geometry would buy, if built}\label{what-the-geometry-would-buy-if-built}}

Supposing both of §3 and §4 could be made rigorous, the payoff is the
thing the whole book has been organized around: \textbf{logical regimes become
regions of \(\mathcal{S}(p)\), not a list of separately-specified formal
systems.} The classical regime is the region where \(\delta^\pm = 0\)
everywhere and \(\Sigma(p) \neq B\) for every \(p\) --- Chapter 22's Glut-Free
Reconstruction condition, read as a subset rather than a proposition.
The paraconsistent regime is the region admitting \(\Sigma(p) = B\) without
requiring \(\delta = 0\) to accompany it --- a strictly larger region,
containing the classical one as a sub-region rather than standing beside
it as an alternative. A relevant-logic-flavored regime would be the
region where \(C(a,p) > 0\) matters for enough actions that propagation
needs the \(\rho_a\)-sensitivity machinery from Chapter 23 to behave
sensibly at all --- not a separate logic bolted on, but a description of
where in \(\mathcal{S}(p)\) an agent's actions live.

If this works, it inverts the usual presentation this book set out to
avoid: instead of ``here is classical logic, and here are its
non-classical rivals,'' the picture becomes one space, with the historical
logics showing up as its more heavily-traveled or more easily-described
regions --- classical logic isn't a rival to paraconsistency, it's the
sub-region of \(\mathcal{S}(p)\) where paraconsistency's extra room happens
not to be needed.

\hypertarget{honest-status}{%
\section{6. Honest status}\label{honest-status}}

Section 1 is proven, in the sense that it restates an already-established
result in geometric language rather than adding new content. Sections 3
and 4 are a program, stated as specifically as this draft can manage
without doing the actual work: two genuine technical questions (does
refinement only ever reveal coherence defects, never manufacture them;
does an action-independent repair distance exist and satisfy a triangle
inequality) stand between ``proposed'' and ``built.'' Section 3 additionally
now carries a settled answer to the question Chapter 15 posed about it:
as proposed, this topology is qualitative-only and doesn't duplicate
dispersion obstruction, which needs scale structure this section doesn't
yet have. Section 5 is the vision the rest of the book has been
implicitly working toward, stated plainly so it can be checked against,
not claimed as delivered. This chapter does what it was supposed to do
on the plan set out earlier: it was generated by the pressure of the
chapters before it, asks questions those chapters make it possible to
ask precisely, and answers only the one Chapter 15 handed it a decision
procedure for.
